
%COVID cancelled standardized testing in the UK and lead to record number of  top grades getting issues for two years straight, and overstreching universities capacities.
How would teachers and schools grade students differently compared to a standardized grading system where teachers have no discretion? The COVID-19 pandemic canceled standardized exams in the UK for two consecutive years, and the final grades were replaced by teachers’ predicted grades. These teacher-assigned grades determined students’ placement outcomes, as the UK higher education system bases admission decisions on applicants’ standardized exam results. The policy led to a 12-percentage-point increase in the number of top grades awarded in 2020, and an even larger increase in 2021, despite the disruptions in learning for the 2021 cohort. Top universities over-accepted students as a majority of their applicants passed the grade requirements, while lower-ranked universities suffered from students shortages. This episode underscores how a change in the grading system reshaped the grade distribution, with far-reaching consequences for the higher education market.

% Why this is economics? Because its localization of information.
I use the complete transfer of grading discretion from a centralized examiner to local teachers to examine the difference across schools in their grading leniency, as well as their grading biases across student backgrounds. To estimate the extent of grading leniency\footnote{I use the term grading leniency as the institutional difference in the students scores in the standardized exam and the teachers assigned grade, while I use grade bias as the difference in the grades assigned across students with different demographic attributes while controlling for their academic ability.} across schools or students, researchers typically compare outcomes under standardized exams with those under teachers’ discretion, aggregate the discrepancies at the school level, and compare them across institutions\footnote{Other papers detect anomalies by comparing students near the cutoff of a passing grade (\cite{neal2008}, \cite{diamond}). Grading differences across demographic groups are estimated by comparing the size of the grading discontinuities across the passing threshold by demographic groups. Such approaches do not compare across exams, and derives grading anomalies using variation across the same test.}. Several studies document systematic differences between teacher-assigned and centrally examined grades (\cite{10.1162/00335530360698441}, \cite{lavy2009}, \cite{hanna2012}, \cite{dee2019}, \cite{Alesina2024}), showing how teachers' extent of grading bias varied by the students gender and race, as well as the attributes related to the teacher or the school. However, no comprehensive comparison of grading leniency across wide range of schools and their underlying causes and consequences exists\footnote{\cite{figlio2006} examines different responses from schools to external incentives to increase student performance, finding how private schools mis-reports their students disability status to hide the bad performance of lower performing students. However they do not focus on grades of the students as the schools in their setting used standardized tests to measure a schools performance.}. The gap reflects the difficulties for researchers to compare across school with a common measure of exam results to be tested against the teachers' assigned grade.

% How schools inflate differently. Static patterns as well as their dynamics.
Using administrative data with final standardized exam  (e.g. \emph{A-levels})  outcomes across each school's cohorts for students that sat the A-levels and students whose A-levels were canceled, I estimate the difference in the grading leniency across schools by comparing the increments in final grades between cohorts. A primary identification threat is that cross-school divergences in grades may reflect pre-existing trends rather than the change in grading method. I show how the exam administrators of the A-levelsbefore the pandemic drew the grade boundaries for the A-level so that the overall grade distribution remains constant overtime after conditioning on the academic performance of the test-taking student population. Moreover, I confirm the stationary property of the grade distribution through multiple stationarity tests, and show that the conditional expected grade of a student follow the same levels overtime. The stationarity property of the expected grade holds even after sub-sampling the sample by school types, students parental income class, and other co-variates, which establishes how the A-levels are designed to balance the difficulty of the A-levels across years by using the test-takers previous academic achievements in a standardized test (e.g.\emph{GCSE}). The annual alignment of student population and the grade distribution allows me to capture the effect of the transfer of grading discretion to local teacher by measuring the deviance between a students' assigned grade to the expected grade that the students would have received if they took the standardized exam. 

My approach for detecting grading leniency across schools resembles the approaches used in \cite{neal2008}; \cite{dee2019}, in which researchers leverage variation in the grading methods used across periods, and estimate the difference in grades between students with academically or demographically similar attributes in different periods\footnote{My approach follows a similar set up as the pandemic completely transferred the grading authority from central examiners to local teachers for the cohorts taking the A-levels in during the pandemic years. }. Additionally, my paper leverages the institutional grading system of the standardized tests to control for any confounders that could potentially influence the final grades. Secondly, I estimate the grading bias across demographic groups using the methods used in grading bias literature (\cite{lavy2018}, \cite{10.1093/qje/qjz008}), which compares the difference in grades between a standardized test and a teacher-assigned grade across demographic groups. I mainly follow the logic of these papers to detect grading biases, by arguing that the composition of students within a classroom was not influenced by the teachers' bias as the switch to teachers predicted grades was completely unpredictable when students were choosing which subjects to take. Lastly, I provide a novel decomposition that identifies whether the extent of grading bias across demographics group varied by the school as well as their driving factors. 

I find substantial heterogeneity in how schools loosened their grading policies, with higher degrees of inflation concentrated in privately funded schools and schools with historically lower average performance in standardized A-level exams. I find a nearly 13 percentage point gap in assigning top grades between schools with quality measures in the bottom 10 decile group and those in the highest 90th decile group. A higher degrees of grade inflation was observed in independent schools across all school quality levels, as their inflation extent was higher than publicly funded schools with similar academic quality. Turning to the channels of grade inflation, I find that lower quality schools inflated their grades in non-quantitative subject courses. This reflects the innate difference between subjects in how much grading discretion are passed to the graders. I find a 20 percentage point difference the bottom 10 decile group schools and the top 90 decile schools in granting top grades in non-quantitative subjects, and little or no evidence across schools in their degree of inflation in quantitative subjects. The inflation levels were largely consistent across the two consecutive years for a majority of the schools, but I observe a downward revisions of grading standards in 2021 at independent schools and schools with stricter grading policy in non-quantitative subjects in the first year.

Turning to grading bias, I find a persistent gender gap in the degrees of grade inflation: conditional on prior performance, female students received top grades by nearly 1.3 and 1.5 percentage point higher chances in 2020 and 2021 respectively. However, the gender bias was concentrated at a subset of schools; I only find a stronger grading biases for female students in schools at higher quality schools in both 2020 and 2021. Furthermore, in 2021 the gender gap increased at schools at extensive margins of the schools quality distribution, and the gender gap remained negligible for schools at the lower end of the quality distribution. Moreover, I find grading bias against ethnic minorities, as white students received top grades with higher chance than any other ethnic groups. Lastly, I find that students from the degrees of grade inflation often positively correlated with the students baseline academic or socio-economic affluence. Such results summarize how teachers did not reverse the expected final grade across students with respect with the students background affluence.

My results contributes to the large empirical literature on grading biases by identifying the institutions that exhibited a stronger grading bias patterns confirmed in the literature\footnote{For example, the grading bias literature finds mixed results in teachers biases towards demographic groups. \cite{10.1093/qje/qjz008} and \cite{lavy2018} finds supporting evidence of gender bias against boys, while \cite{hanna2012} and \cite{diamond} do not find any evidence.}. Firstly, I show how grade inflation was more observed in secondary schools with lower average grades, presumably drive by differences in student compositions as the academic ability of the student upward bounds the extent of the grade inflation\footnote{The results are similar to \cite{figlio2004}, \cite{10.1162/00335530360698441}, and \cite{dee2019} findings of stronger concentration of grade manipulation at lower quality schools.}. The upper bound of the grade distribution caps the extent of grade inflation for students with better economic affluence (\cite{hanna2012}, \cite{burgess2013}), especially at school with higher quality. I only observe the positive relationship between parental income and the extent of inflation in higher quality schools. Secondly, I discover a novel channel (e.g. \emph{Subject content}) of grading discretion across students inside a school. I find subject content sets differential degrees of grading discretion onto teachers and amplifies the existing biases between demographic groups. Also the grading standards within non-quantitative subjects were fluctuate, as schools with stricter grading standards in the first year swiftly revised them downwards in the second year. The findings confirm the documents the grade inflation patterns in \cite{denning2023} in which average GPA at universities increases, and also provides the channels that schools exploited to improve their students placements. 

My results also contributes on the effect of admission policies on university diversity. Test optional policies remove the requirements for students to submit a standardized test, with advocates citing how test-optional policies allows universities to allow a broader range of students for admissions (\cite{10.1257/aer.20231407}). Recent empirical work point to the inefficiency of such policies, citing enrollment reduction(\cite{NBERw34260}) and potential crowding out of good grades from  (\cite{NBERw33389}) of disadvantaged students. My paper shows how the removal of standardized grade in the UK evened the grade outcome across students with contrasting backgrounds, but the improvements in grades did not translate into better placements of disadvantaged students. My finding highlight how removal of standardized grade levels up the signals that disadvantaged students can capitalized on, but with limited impact on diversifying the actual pool of applying candidates that selective university can choose from. 

The paper is organized as follows: Section 2 provides the institutional background of the higher education system in the UK as well as the details of the policy set in place during the pandemic. Section 3 describes the data source I use in the analysis. Section 4 presents the empirical model and discuss the validity of the estimates along with the institutional details. Section 5 presents the results on heterogeneity of grading patterns across schools, as well as student level patterns, the dynamics of grading patterns, as well as their influence in placing students to their top choices. Section 6 discuss the magnitudes of the results with the literature on grading biases. Section 7 concludes with policy discussion on how my finding can inform policy makers with thinking about test-optional policies. 